\chapter{Introducción, motivación y objetivos}
\label{cap:introduccion}

El presente capítulo introduce el contexto en el que se enmarca este
Trabajo Fin de Máster, en adelante TFM, describe la motivación que
justifica su desarrollo y delimita los objetivos que guían la solución
propuesta. Para ello, en primer lugar se sitúa el trabajo con respecto al
Trabajo Fin de Grado, en adelante TFG, previo del mismo autor, del que
parte directamente. A continuación, se expone el problema de partida y la
motivación que lo sustenta, se establecen los objetivos general y
específicos, se recogen los aspectos éticos y de privacidad que afectan al
tratamiento de los datos empleados y se relacionan los objetivos con los
resultados de aprendizaje del máster. Finalmente, se presenta la
estructura seguida por el resto del documento.

\section{Contexto y continuidad con el Trabajo Fin de Grado}
\label{sec:introduccion_contexto}

Este TFM se construye directamente sobre el trabajo ya cerrado y defendido
en el TFG del mismo autor~\cite{tfg_open_cvn}, del que hereda el punto de
partida técnico: una capa estructural reproducible generada a partir de
los esquemas oficiales de CVN, una normalización trazable de sus
metadatos, y un formato abierto, validable y extensible, Open CVN JSON,
con su propio contrato de parseo y validación.

El presente trabajo reutiliza ese resultado como una de sus dos fuentes de
datos, no como objeto de estudio: se reutilizan directamente el esquema
JSON Schema de Open CVN, su contrato de validación ya implementado, y las
convenciones de trazabilidad de procedencia, en inglés
\textit{provenance}, establecidas en la capa semántica del TFG. Por ello,
este documento no vuelve a definir qué es CVN ni qué es Open CVN,
remitiendo esa exposición al TFG anterior, y se limita a describir cómo
dicho formato, ya cerrado, se ingiere, valida y procesa dentro de la
plataforma distribuida que aquí se construye.

Este cambio de perspectiva, de un formato y una herramienta de un único
usuario a una plataforma que ingiere, fusiona y procesa datos a escala,
responde también a un cambio de programa académico. Mientras que el TFG se
enmarcaba en la intensificación de Computación, este TFM pertenece al
Máster en Big Data y Computación en la Nube, cuyos resultados de
aprendizaje se detallan en la Sección~\ref{sec:introduccion_outcomes}.

\section{Problema de partida y motivación}
\label{sec:introduccion_motivacion}

La información curricular y de trayectoria investigadora se encuentra
dispersa en fuentes heterogéneas, de distinta naturaleza y de distinto
volumen. Por un lado, documentos curriculares como CVN describen la
actividad de un investigador de forma estructurada pero unitaria, pensada
para su presentación individual ante convocatorias, evaluaciones o
acreditaciones. Por otro lado, registros públicos como los perfiles de
ORCID, siglas en inglés de \textit{Open Researcher and Contributor
ID}~\cite{orcid_about}, agregan identificadores, filiaciones y producción
científica de investigadores de todo el mundo a escala global. Solo el
fichero público de datos de ORCID de 2025 reúne más de 26 millones de
registros~\cite{orcid_public_data_file_2025}.

Los sistemas de información de investigación, conocidos por sus siglas en
inglés como CRIS, de \textit{Current Research Information Systems},
interoperan datos de este tipo entre instituciones mediante modelos
comunes como CERIF~\cite{eurocris_cerif} o plataformas como
VIVO~\cite{vivo_project}. Sin embargo, lo hacen sobre infraestructura
institucional cerrada, sin publicar una plataforma de referencia abierta,
autoalojada y reproducible que combine específicamente currículos
normalizados como CVN con perfiles públicos como ORCID. Estas dos
naturalezas de dato, documental y de registro masivo público, rara vez se
integran, procesan conjuntamente y publican como indicadores agregados
dentro de una misma plataforma de este tipo.

Esta ausencia responde a un problema de arquitectura. Integrar fuentes de
esta naturaleza exige:

\begin{itemize}
  \item \textbf{Almacenamiento escalable}, capaz de crecer con el volumen
  real de ORCID.
  \item \textbf{Un formato de tabla que preserve evolución de esquema y
  trazabilidad}, dado que CVN y ORCID no comparten esquema.
  \item \textbf{Procesamiento distribuido tolerante a fallos}, para
  transformar ese volumen en un tiempo razonable.
  \item \textbf{Orquestación de flujos de trabajo}, que encadene ingesta,
  transformación y publicación de forma reproducible.
  \item \textbf{Una capa de resolución de identidad}, capaz de fusionar
  registros de fuentes distintas que describen a la misma persona.
\end{itemize}

Ninguno de estos elementos formaba parte del alcance del TFG, centrado
deliberadamente en un único documento curricular.

La motivación de este trabajo consiste, por tanto, en demostrar, mediante
una plataforma real desplegada y no mediante un diseño exclusivamente
documental, las competencias propias de un perfil de Big Data y
Computación en la Nube:

\begin{itemize}
  \item Levantar infraestructura distribuida.
  \item Ingerir y fusionar fuentes heterogéneas.
  \item Procesar los datos a escala con tolerancia a fallos.
  \item Evaluar el sistema con datos y cifras reales obtenidos de su
  propia ejecución, en lugar de resultados simulados o proyectados.
\end{itemize}

\section{Objetivos}
\label{sec:introduccion_objetivos}

El objetivo general de este Trabajo Fin de Máster es \textit{diseñar,
desplegar y evaluar una plataforma que integre datos curriculares de CVN
sintético y datos de investigación de ORCID real, los procese de forma
distribuida con resolución de identidad entre fuentes, y publique un
conjunto reducido de indicadores de investigación en un panel de
visualización, evaluando el sistema mediante un \textit{benchmark} de
rendimiento acotado}. Los objetivos que siguen se formulan de forma
agnóstica respecto a la arquitectura y las tecnologías concretas que los
satisfacen. Esa elección se presenta y se justifica frente a sus
alternativas en el Capítulo~\ref{cap:estado_arte}.

Para alcanzar esta meta global se han definido los siguientes objetivos
específicos:

\begin{enumerate}
  \item \textbf{Desplegar la infraestructura base de la plataforma.}
  Comprende la puesta en marcha del entorno de cómputo y de los servicios
  de almacenamiento, persistencia y orquestación sobre los que se apoya el
  resto del sistema.

  \item \textbf{Validar el mecanismo central de almacenamiento y
  procesamiento distribuido.} Antes de construir cualquier flujo de datos
  de negocio sobre la plataforma, se comprueba de forma aislada que dicho
  mecanismo funciona de extremo a extremo.

  \item \textbf{Incorporar datos ORCID reales a la plataforma.} Combina el
  enriquecimiento puntual de registros concretos con la incorporación de un
  volumen real de datos, ambos obtenidos de fuentes públicas de ORCID.

  \item \textbf{Incorporar currículos CVN a la plataforma, validados frente
  al formato Open CVN.} Estos currículos se obtienen por generación
  sintética a partir de datos públicos, no por recolección de currículos
  reales, tal y como se justifica en la Sección~\ref{sec:introduccion_etica}.

  \item \textbf{Integrar los datos de ORCID con los datos CVN por medio del
  formato Open CVN, resolviendo la identidad de los registros que
  corresponden a la misma persona.} Constituye la fusión efectiva entre las
  dos fuentes de datos del sistema.

  \item \textbf{Calcular indicadores de investigación agregados y
  publicarlos de forma consistente.} Los indicadores deben quedar
  disponibles para su consumo posterior sin estados intermedios ni
  publicaciones parciales.

  \item \textbf{Visualizar los indicadores publicados en un panel
  accesible.} Permite comprobar que los indicadores calculados son
  consumibles más allá del propio sistema que los produce.

  \item \textbf{Medir el efecto del grado de paralelismo del procesamiento
  distribuido sobre el tiempo de ejecución y la fiabilidad, a distintas
  escalas de datos.} Constituye el \textit{benchmark} de rendimiento
  acotado exigido por el planteamiento del proyecto.

  \item \textbf{Encontrar y corregir defectos de robustez operativa mediante
  el uso real del sistema completo.} Este objetivo, sin equivalente directo
  en el TFG anterior, exige documentar una guía de reproducción del entorno
  desde cero y registrar los fallos de infraestructura detectados y
  corregidos como parte de la evaluación del sistema.
\end{enumerate}

La Tabla~\ref{tab:tfm_objetivos_capitulos} relaciona estos objetivos con
los capítulos principales en los que se desarrollan. Es en esos capítulos,
no aquí, donde cada objetivo se concreta en una arquitectura y unas
tecnologías determinadas.

\begin{table}[htbp]
  \centering
  \renewcommand{\arraystretch}{1.2}
  \begin{tabular}{p{0.12\textwidth} p{0.58\textwidth} p{0.18\textwidth}}
    \hline
    \textbf{Objetivo} & \textbf{Descripción resumida} & \textbf{Capítulos} \\
    \hline
    OE1 & Despliegue de la infraestructura base de la plataforma & 3 \\
    OE2 & Validación del mecanismo central de almacenamiento y procesamiento & 3 \\
    OE3 & Incorporación de datos ORCID reales & 4 \\
    OE4 & Incorporación de currículos CVN, validados frente a Open CVN & 4 \\
    OE5 & Integración de ORCID con CVN por medio de Open CVN & 5 \\
    OE6 & Cálculo y publicación consistente de indicadores & 5 \\
    OE7 & Visualización de los indicadores en un panel accesible & 6 \\
    OE8 & Benchmark de rendimiento por grado de paralelismo y escala de datos & 6 \\
    OE9 & Robustez operativa mediante uso real y guía de reproducción & 6 \\
    \hline
  \end{tabular}
  \caption{Relación entre objetivos específicos y capítulos del documento.}
  \label{tab:tfm_objetivos_capitulos}
\end{table}

\section{Aspectos éticos y de privacidad en el tratamiento de datos}
\label{sec:introduccion_etica}

El tratamiento de datos personales de investigadores exige distinguir
entre las dos fuentes empleadas. Los perfiles de ORCID se usan como dato
real porque su naturaleza es, por diseño, pública: un identificador ORCID
existe precisamente para ser consultado de forma abierta e interoperable
entre sistemas de gestión de la investigación, así que su uso aquí no
introduce ningún tratamiento distinto del que la propia infraestructura de
ORCID habilita.

Los documentos CVN, en cambio, no disfrutan de esa misma naturaleza
pública: un currículo completo es información personal habitualmente
empleada en solicitudes o evaluaciones privadas, sin un corpus público
legítimo del que obtener volumen. Recolectarlos supondría tratar datos
personales de terceros sin su consentimiento, un problema de privacidad
real para un repositorio público y un trabajo defendido públicamente. Por
ello, este documento descarta recolectar CVN reales a volumen y opta por
generar currículos sintéticos, válidos frente al esquema oficial y
sembrados con campos públicos reales de ORCID, como nombres, obras,
filiaciones y fechas, de forma que resulten realistas sin representar a
ninguna persona real. Esa misma siembra crea el vínculo por identificador
ORCID que la resolución de identidad entre fuentes necesita como clave de
fusión.

\section{Resultados de aprendizaje del máster}
\label{sec:introduccion_outcomes}

El desarrollo de este proyecto se alinea con los resultados de aprendizaje
del Máster en Big Data y Computación en la Nube, y no con las
competencias de la intensificación de Computación bajo las que se evaluó
el TFG anterior. La Tabla~\ref{tab:tfm_learning_outcomes} recoge los seis
resultados de aprendizaje sobre los que este trabajo aporta evidencia
concreta, junto con los capítulos en los que dicha evidencia se desarrolla.

\begin{table}[htbp]
  \centering
  \renewcommand{\arraystretch}{1.2}
  \begin{tabular}{p{0.1\textwidth} p{0.58\textwidth} p{0.2\textwidth}}
    \hline
    \textbf{Resultado} & \textbf{Descripción} & \textbf{Capítulos} \\
    \hline
    CN02 & Arquitecturas de tratamiento masivo, almacenamiento y \textit{pipelines} & 2, 3 \\
    HA01 & Adquisición, fusión y análisis de múltiples fuentes & 4, 5 \\
    HA02 & Procesamiento distribuido y optimización de recursos & 5, 6 \\
    HA03 & Orquestación ETL y almacenamiento escalable & 3, 5 \\
    CP01 & Planificación y despliegue de una solución Big Data & 2, 3, 6 \\
    CP04 & Proyecto profesional original de Big Data & 6, 7 \\
    \hline
  \end{tabular}
  \caption{Relación entre resultados de aprendizaje del máster y capítulos del documento.}
  \label{tab:tfm_learning_outcomes}
\end{table}

Ninguno de estos resultados se demuestra de forma primaria en el presente
capítulo introductorio. Su cumplimiento efectivo se retomará, con
evidencia concreta extraída del sistema desarrollado, en el
Capítulo~\ref{cap:conclusiones}.

\section{Estructura del documento}
\label{sec:introduccion_estructura}

El resto del documento se organiza en seis capítulos adicionales, desde la
justificación tecnológica hasta la evaluación y las conclusiones:

\begin{itemize}
  \item \textbf{Capítulo 2: Estado del arte y decisiones arquitectónicas.}
  Alternativas tecnológicas consideradas para cada componente de la
  plataforma y justificación de las decisiones adoptadas.

  \item \textbf{Capítulo 3: Infraestructura: despliegue del clúster y
  servicios base.} Despliegue del clúster Kubernetes local y de los
  servicios sobre los que se apoya el resto del sistema.

  \item \textbf{Capítulo 4: Ingesta y fusión de fuentes heterogéneas.}
  Ingesta de datos ORCID reales, generación de CVN sintéticos y su
  aterrizaje con procedencia en la capa \textit{bronze}.

  \item \textbf{Capítulo 5: Procesamiento distribuido: transformación,
  resolución de entidades e indicadores.} Validación, normalización y
  resolución de identidad entre fuentes, y cálculo y publicación de
  indicadores agregados.

  \item \textbf{Capítulo 6: Visualización, evaluación de rendimiento y
  endurecimiento.} Panel de indicadores, \textit{benchmark} de rendimiento
  y defectos de robustez encontrados y corregidos mediante uso real.

  \item \textbf{Capítulo 7: Conclusiones, competencias y trabajo futuro.}
  Cumplimiento de objetivos, resultados de aprendizaje demostrados,
  limitaciones y líneas futuras.
\end{itemize}
